% Reconstructed 2026-11-20 from complete authored proof chunks and audited formulas.
% This is a new source identity, not a byte-identical recovery of the lost 13,054-character source.
\Needspace{10\baselineskip}
\subsection{A.154. Spatial Lipschitz factors and real higher surjectivity}

\begin{saresult}{A.154}{Proved under stated hypotheses}{Finite spatial regularization and higher comparison over the reals}
Let $M$ be an o-minimal real closed field, let $X$ be a compact
finite PL polyhedron, and fix finitely many marked subpolyhedra.
Let $P=[0,1]^k$ and let $U_p:X\to X$ be a jointly continuous
definable family of homeomorphisms preserving the marks and
equal to the identity for $p\in\partial P$. There are jointly
continuous definable families of homeomorphisms $r_p,s_p$,
preserving the marks and equal to the identity on $\partial P$,
such that
\[
 U_p=s_pr_p^{-1},
 \tag{154.1}
\]
and both families have one common uniform spatial Lipschitz
bound. If $U$ fixes a marked subpolyhedron pointwise, the two
factors agree there. The construction uses any fixed finite
target dimension.

If $M$ is an o-minimal expansion of the real field and $B$ is
a closed compact PL manifold, then
\[
 \pi_kP_{\rm PL}(B)\longrightarrow\pi_kP_{\rm def}(B)
 \quad\hbox{is surjective for every }k\ge1.
 \tag{154.2}
\]
No injectivity assertion follows from this theorem alone.
\end{saresult}

\begin{proof}
Use maximum norms.

\textbf{Finite variation and its stability.}\quad
For a definable family of continuous scalar functions on
$[0,1]$, o-minimal monotonicity gives an ordinary finite
uniform bound $N$ for the number of monotonicity intervals.
For each member $f$ and $0\le t\le1$ its variation is exactly
\[
 V_f(t)=
 \max_{0=t_0\le t_1\le\cdots\le t_N=t}
       \sum_{\nu=1}^N|f(t_\nu)-f(t_{\nu-1})|.
 \tag{154.3}
\]
Repeated partition points are allowed. A partition by the
monotonicity intervals realizes the ordinary variation, so
the displayed finite maximum is exact. It is definable,
and its maximum exists by definable compactness. Comparison
on the same partition gives
\[
 |V_f(t)-V_g(t)|\le2N\|f-g\|_\infty.
 \tag{154.4}
\]
For a finite vector of $d$ scalar functions, sum their
variations; the comparison constant is then $2dN$.
Variation is additive on adjacent intervals, and hence
\[
 \|f(v)-f(u)\|_\infty\le V_f(v)-V_f(u)
 \quad(0\le u\le v\le1).
 \tag{154.5}
\]
The variation depends jointly continuously on the parameters
and on $t$. To check the endpoint variable, apply (154.4)
to the rescaled lines $z\mapsto f(tz)$ on the fixed unit
interval. Definable uniform continuity on a compact product
gives the required uniform convergence of these lines.
These arguments work over every o-minimal real closed field.

\textbf{The mixed coordinate-clock induction.}\quad
Let $C=[0,1]^n$, and let
$E:C\times P\to[0,1]^a$ be any continuous definable vector
map. Begin with
\[
 \chi^0_p=\mathrm{id}_C,\qquad
 F^0(x,p)=(x,E(x,p)).
\]
At stage $i$, suppose the vector
$F^{i-1}=(\chi^{i-1},E\chi^{i-1})$ is uniformly Lipschitz
in every already processed coordinate $x_j$, $j<i$.
Fix the other coordinates and $p$ as a profile $b$, and
write $f_b(t)$ for its $i$th coordinate line. Define
\[
 A_b(t)=t+V_{f_b}(t),\qquad
 D_b=A_b(1),\qquad
 \xi_b(u)=A_b^{-1}(D_bu).
 \tag{154.6}
\]
The function $A_b$ is continuous and strictly increasing;
$A_b(v)-A_b(u)\ge v-u$. Its inverse exists and depends
jointly continuously on the profile. Replace $x_i$ by
$\xi_b(x_i)$, leaving all other coordinates and $p$ fixed,
to obtain a parameter-preserving homeomorphism $\Psi_i$.
Its inverse changes that coordinate to $A_b(x_i)/D_b$.
Both maps preserve each coordinate face. Set
\[
 \chi^i=\chi^{i-1}\Psi_i,\qquad
 F^i=F^{i-1}\Psi_i.
\]

Here is the needed mixed estimate. Let two profiles $b,c$
differ in a previously processed coordinate by $\delta$,
with old Lipschitz bound $K$. If the vector has $d$
components and its line monotonicity bound is $N$, put
$C_0=2dN$. Formula (154.4) gives
\[
 \sup_t|A_b(t)-A_c(t)|\le C_0K\delta,\qquad
 |D_b-D_c|\le C_0K\delta.
\]
For $h=\xi_b(u)$ and $h'=\xi_c(u)$,
\[
 |A_c(h)-A_c(h')|\le2C_0K\delta.
\]
Using (154.5) to control the change along the old line,
rather than assuming that line was Lipschitz, yields
\[
 \|f_b(h)-f_c(h')\|_\infty
       \le(1+2C_0)K\delta.
 \tag{154.7}
\]
For a change in the new coordinate $u$,
\[
 \|f_b(\xi_b(u))-f_b(\xi_b(u'))\|_\infty
       \le D_b|u-u'|.
 \tag{154.8}
\]
The heights $D_b$ have a uniform bound by finite variation
and bounded coordinate ranges. Thus the induction preserves
all earlier coordinate estimates and adds the new one.
After $n$ stages, changing coordinates one at a time proves
uniform global spatial Lipschitz bounds for both $\chi$
and $E\chi$. The evolving domain map is included in the
vector throughout; omitting it would not prove a bound
for the first factor.

If all vector components have standard bounded ranges,
these final constants have standard integer bounds.
Each line variation is bounded by its range times the
ordinary finite monotonicity count. Equations (154.7)--(154.8)
then involve only finitely many ordinary numerical operations.
In particular this conclusion holds on a fixed standard
unit-cube model. No integer is chosen larger than an
arbitrary infinite field element.

If $E_p=\mathrm{id}_C$ for $p\in\partial P$ and $a=n$,
then on this boundary the evolving vector is $(x,x)$.
Each processed line has variation $2t$, so $A_b(t)=3t$,
$D_b=3$, and $\xi_b(u)=u$. Hence every clock fixes the
parameter boundary pointwise.

\textbf{A PL cubical model with all marks.}\quad
Triangulate $X$ compatibly with the finite marks and number
the vertices by $1,\ldots,v$. For the standard realization
of the abstract complex $K$, set
\[
 Q=\{z\in[0,1]^v:
       \operatorname{supp}(z)\in K,\ \max_i z_i=1\}.
 \tag{154.9}
\]
It is the finite union of the coordinate faces
$z_i=1$, $z_j=0$ for $j\notin\sigma$, where
$\sigma\in K$ and $i\in\sigma$. A marked subcomplex is
the corresponding union of such faces.

There is a PL homeomorphism $b:X\to Q$ preserving these
marks. On the barycentric subdivision of the standard
realization, send the barycentre of $\sigma$ to its
characteristic vector $\mathbf1_\sigma$ and extend
affinely over chains. These chains give the usual
monotone-coordinate triangulations of the displayed
coordinate faces, and agree on their intersections.
They triangulate all of $Q$. Thus this is a finite
PL homeomorphism. One may check the underlying topology
also by $z\mapsto z/\sum z_i$ and
$\lambda\mapsto\lambda/\max\lambda_i$, but those rational
maps are not the conjugacy used here. The actual map $b$
and its inverse are Lipschitz on their finite polyhedra.

\textbf{Extending a vector map, not homeomorphisms.}\quad
Conjugate $U$ to $U^Q$ on $Q$. On the closed polyhedral
subcomplex
\[
 Z=(Q\times P)\cup(C\times\partial P),
 \qquad C=[0,1]^v,
\]
prescribe $U^Q$ on the first piece and the identity on the
second. The prescriptions agree and define a continuous
definable map $E_0:Z\to C$.

Take a finite compatible triangulation of $C\times P$ and
barycentrically subdivide so that $Z$ is full. On a simplex,
let $\alpha(w)$ be the sum of its barycentric weights at
vertices lying in $Z$. If $\alpha>0$, let $q(w)$ be the
normalized barycentric sum of those vertices. Fullness
ensures $q(w)\in Z$. For a fixed $c\in C$, put
\[
 E(w)=\alpha(w)E_0(q(w))+(1-\alpha(w))c
      \quad(\alpha(w)>0),\qquad E(w)=c\quad(\alpha(w)=0).
 \tag{154.10}
\]
The formulas agree on shared faces and extend $E_0$.
They are continuous at $\alpha=0$ because $E_0$ is bounded;
their values lie in the convex cube $C$.

Apply the coordinate clocks to this vector map. They
preserve every coordinate face, hence $Q$ and every mark.
Their restrictions to $Q$ are homeomorphisms. Define
\[
 r_p=b^{-1}\chi_pb,\qquad s_p=U_pr_p.
 \tag{154.11}
\]
Since $E=U^Q$ on $Q$, the spatial bounds for $\chi$ and
$E\chi$ give the required bounds for $r,s$ after fixed
PL conjugation. They satisfy (154.1), preserve all marks,
and are identity on the parameter boundary. If $U$ fixes
a mark pointwise, preservation by $r$ implies $s=r$ on it.
The total maps and their inverses are jointly continuous;
this also follows from the continuous bijection of compact
definable products. This proves the first assertion.

\textbf{From spatial factors to parameter factors over the reals.}\quad
Now take $X=N=B\times I$ over $\mathbb R$, with its bottom
and top marked. The factors agree on the bottom. Separate
their graphs by a fixed tag larger than the diameter of
$N\times N$:
\[
 \mathcal G_p=(\{0\}\times\Gamma r_p)
                   \cup(\{T\}\times\Gamma s_p).
 \tag{154.12}
\]
The Hausdorff distance of these unions is the maximum of
the two individual graph distances, since any same-tag
point is closer than every opposite-tag point.
This definable compact-set family is continuous and
constant on the parameter boundary.

Apply the exact relative cube assertion A.152, using the
fixed affine identification $[0,1]^k\cong[-1,1]^k$ if that
cube convention is used there. It supplies a definable
Lipschitz map $g:P\to P$, equal to the identity on
$\partial P$, for which $\mathcal G\circ g$ is Hausdorff
Lipschitz. Let $L$ bound each pulled individual graph
family and let $K$ bound both factors spatially.
For each $x$, graph matching supplies an $x'$ with
\[
 \|x-x'\|_\infty\le L\|p-q\|_\infty,\qquad
 \|r_{g(p)}(x)-r_{g(q)}(x')\|_\infty
                           \le L\|p-q\|_\infty.
\]
Consequently
\[
 \|r_{g(p)}(x)-r_{g(q)}(x)\|_\infty
                \le(1+K)L\|p-q\|_\infty,
 \tag{154.13}
\]
and the same holds for $s$. Thus $r\circ g,s\circ g$ have
uniform forward parameter bounds. No inverse bound is used.
They still agree on the bottom and are both the identity
on the parameter boundary. The pair criterion A.148 gives
a jointly PL representative of their quotient $U\circ g$.

Finally the homotopy
$U_{(1-t)p+t g(p)}$ is based, bottom-fixed, and jointly
continuous definable. It connects $U$ to $U\circ g$,
proving (154.2). Empty bases are immediate, and the marked
construction treats disconnected and zero-dimensional
bases componentwise.
\end{proof}

